\documentclass[11pt]{article}
\usepackage[a4paper,margin=25mm]{geometry}
\usepackage[T1]{fontenc}
\usepackage{lmodern,microtype,amsmath,amssymb,amsthm,mathtools}
\usepackage{booktabs,array,enumitem,xcolor}
\usepackage[colorlinks=true,linkcolor=blue!55!black,urlcolor=blue!55!black,citecolor=blue!55!black]{hyperref}
\hypersetup{pdftitle={No universal finite collapse of the upper PPT entanglement-cost hierarchy},pdfauthor={Dev D. Goyal},pdfsubject={Mathematical proof, Lean 4 verification, and qualified novelty assessment}}
\usepackage{fancyhdr}
\setlength{\headheight}{14pt}
\pagestyle{fancy}\fancyhf{}
\fancyhead[L]{Noncollapse of the upper PPT hierarchy}
\fancyhead[R]{30 October 2026}
\fancyfoot[C]{\thepage}
\setlength{\parskip}{5pt}\setlength{\parindent}{0pt}
\setlist{itemsep=2pt,topsep=4pt}
\newtheorem{theorem}{Theorem}
\newtheorem{lemma}{Lemma}
\newtheorem{proposition}{Proposition}
\newcommand{\Tr}{\operatorname{Tr}}
\newcommand{\C}{\mathbb C}
\newcommand{\R}{\mathbb R}
\newcommand{\one}{\mathbf1}
\newcommand{\syn}{\mathcal R}
\newcommand{\E}{\mathbb E}
\newcommand{\diag}{\operatorname{diag}}
\title{No universal finite collapse of the\\upper PPT entanglement-cost hierarchy}
\author{Dev D. Goyal\thanks{Fully autonomous agentic research; see the disclosure at the end. No expert validation or external peer review is claimed.}}
\date{30 October 2026}

\begin{document}
\maketitle
\begin{abstract}
For every integer $q\ge1$, there is a bipartite density matrix $\rho_q$ for which
$\kappa_{q+1}(\rho_q)<\kappa_q(\rho_q)$ in the original upper hierarchy of
Lami, Mele, and Regula. The witnesses can be chosen as full-rank product-Bell
diagonal states on $\C^{d_q}\otimes\C^{d_q}$, where
$d_q=2^{(2q+4)4^q}$. Consequently, no fixed finite upper level is exact for
all bipartite states.

The proof constructs arbitrarily long strictly positive reverse orbits of
absolute partial transposition. A finite weighted-sign argument extends an
orbit while controlling every coordinate. Strict positivity then turns a
known telescoping dual identity into a quantitative obstruction for every
shorter feasible upper chain. Bell expectations transfer the obstruction to
unrestricted complex auxiliary matrices. The central all-depth theorem,
including its physical state assumptions and the original SDP constraints,
is formalized in Lean 4 with pinned dependencies.

The located primary literature already proves arbitrary-depth noncollapse of
the lower $\chi$ hierarchy and explicitly leaves this upper $\kappa$ question
open. The originality assessment concerns the positive inverse-orbit
construction and the resulting all-depth upper separation. It is supported
by targeted searches before and after the proof, but historical priority
remains qualified. The dimension bound is deliberately generous; the result
does not claim efficient small witnesses, noncollapse in one fixed dimension,
or impossibility of other single-letter formulas.
\end{abstract}
\paragraph{Verification status.} The all-depth theorem and the full-rank state statement were replayed from source for this release with the Lean kernel at trust zero. All 17 mathematical modules and the complete 256-declaration axiom audit succeeded, and every audited declaration depends only on \texttt{propext}, \texttt{Classical.choice} and \texttt{Quot.sound}.

\clearpage
\tableofcontents
\newpage

\section{The theorem and its precise scope}
For partial transposition $\Gamma$ on the second local system, define
\begin{equation}\label{eq:kappa}
\kappa_q(\rho)=\inf\left\{\Tr S_{q-1}:\begin{array}{l}
S_{-1}=\rho,\quad S_i\succeq\pm S_{i-1}^{\Gamma}\quad(0\le i<q),\\
S_{q-1}^{\Gamma}\succeq0
\end{array}\right\}.
\end{equation}
This is the original indexing and terminal constraint in
Lami--Mele--Regula~\cite{LMR}. The auxiliary matrices range over the entire
complex matrix space. The PSD sandwich constraints force Hermiticity and
positivity of the auxiliaries, so their traces are real. The formal definition
uses the real part of the trace and proves this correspondence.

\begin{theorem}[Strictness at every adjacent upper level]\label{thm:main}
For every $q\ge1$, put
\[
n_q=(2q+4)4^q,\qquad d_q=2^{n_q}.
\]
There is a positive-definite, trace-one matrix $\rho_q$ on
$\C^{d_q}\otimes\C^{d_q}$, diagonal in a product-Bell basis with strictly
positive coefficients, such that
\begin{equation}\label{eq:main}
\kappa_{q+1}(\rho_q)<\kappa_q(\rho_q).
\end{equation}
Thus there is no integer $q\ge1$ for which $\kappa_q=\kappa_{q+1}$ holds
on every finite-dimensional bipartite quantum state.
\end{theorem}

The construction is an existence proof over a finite collection of sign
vectors, with an explicit dimension bound. It does not provide a short
closed formula for the successful signs or a practical table of all state
coefficients. This distinction does not affect the statement in
Theorem~\ref{thm:main}: Lean proves that the required state exists, rather
than accepting a missing numerical certificate as an assumption.

The operational consequence uses the published bound
\begin{equation}\label{eq:operational}
E^{\rm exact}_{c,\rm PPT}(\rho)\le\log_2\kappa_r(\rho)
\qquad(r\ge1).
\end{equation}
For the witness at level $q$, this implies
$E^{\rm exact}_{c,\rm PPT}(\rho_q)<\log_2\kappa_q(\rho_q)$.
Therefore no fixed upper level gives a universal exact PPT-cost formula.
Equation~\eqref{eq:operational} is established prior mathematics
\cite[Theorem S26, Eq.\ (S89)]{LMR}; the operational theory is not introduced
as a Lean axiom or claimed to be re-formalized in this project.

A further consequence uses the published regularization identity
$E_{\kappa,q}^{\infty}=E^{\rm exact}_{c,\rm PPT}$
\cite[Remark S27, Eq.\ (S98)]{LMR}. For each finite $q$, the witness must
satisfy $\kappa_q(\rho_q^{\otimes m})<\kappa_q(\rho_q)^m$ for at least one
integer $m\ge2$; otherwise its regularized logarithmic value could not be
smaller than its one-copy value. Thus every finite upper level fails tensor additivity
somewhere. This is a consequence of the formally proved separation and a
cited prior theorem, not a separately formalized operational result.

The quantifiers matter. The dimension is allowed to grow with $q$. This
does not prove that one fixed state or one fixed dimension has infinitely
many strict upper improvements. The construction itself reaches a simple
terminal point after $q+1$ steps. The theorem also does not rule out a
different single-letter characterization of exact PPT cost.

\section{Bell coordinates and the unrestricted operator problem}
Use the one-pair Bell order $(\Phi^+,\Phi^-,\Psi^+,\Psi^-)$ and the matrix
\begin{equation}\label{eq:T}
T_1=\frac12\begin{pmatrix}
1&1&1&-1\\1&1&-1&1\\1&-1&1&1\\-1&1&1&1
\end{pmatrix},\qquad T=T_1^{\otimes n}.
\end{equation}
Let $d=2^n$ and $N=4^n=d^2$. The matrix $T$ is symmetric, orthogonal,
involutive, fixes $\one$, and has entries of absolute value $1/d$.
In particular,
\begin{equation}\label{eq:flat}
T^2=I,\quad T\one=\one,\quad |T_{ij}|=d^{-1},\quad
\sum_i(Tx)_i^2=\sum_i x_i^2.
\end{equation}

Let $P_j$ be the rank-one product-Bell projectors and write
\[
\syn(x)=\sum_{j=1}^N x_jP_j.
\]
The Bell bridge proves
\begin{equation}\label{eq:bridge}
\syn(x)^\Gamma=\syn(Tx),\quad
\Tr\syn(x)=\sum_jx_j,\quad
x\ge0\Longrightarrow\syn(x)\succeq0.
\end{equation}
Strictly positive coefficients give a positive-definite synthesized matrix.

For an arbitrary complex matrix $X$, define its Bell expectations by
\[
\mathcal E(X)_j=\operatorname{Re}\Tr(P_jX).
\]
The formal bridge establishes
\begin{equation}\label{eq:expect}
X\succeq0\Longrightarrow\mathcal E(X)\ge0,\qquad
\mathcal E(X^\Gamma)=T\mathcal E(X),\qquad
\sum_j\mathcal E(X)_j=\operatorname{Re}\Tr X.
\end{equation}
Also $\mathcal E(\syn(x))=x$.

Consequently, every feasible complex-matrix chain for $\syn(x)$ yields a
coefficient chain
\begin{equation}\label{eq:chain}
s_0=x,\qquad s_{i+1}\ge|Ts_i|\ (0\le i<q),\qquad Ts_q\ge0,
\end{equation}
with the same real trace objective $\sum_js_q(j)$. Conversely, every such
coefficient chain synthesizes a feasible matrix chain. Lower bounds below
therefore apply to all complex auxiliaries. No assertion that an optimizer
is Bell diagonal is needed.

The manuscript uses the convenient chain indexing $s_0=x$ and $s_q$ terminal.
It matches \eqref{eq:kappa} by $s_i=\mathcal E(S_{i-1})$. The Lean definitions
and compatibility lemmas make this shift explicit.

\section{A positive reverse-step lemma}
Define the absolute-transform map
\[
F(x)=|Tx|,
\]
with absolute value coordinatewise. The central new construction extends a
reverse orbit $F(w)=v$ while preserving strict positivity of $w$.

\begin{lemma}[Finite weighted-sign inverse]\label{lem:inverse}
Suppose $0\le\varepsilon<1$,
\[
|v_i-1|\le\varepsilon\quad\hbox{for every }i,
\qquad \sum_i v_i^2=N.
\]
Let $a,t>0$ satisfy
\begin{equation}\label{eq:gates}
\frac{2t}{d}\le1,\qquad
2N\exp(2\varepsilon t^2)<\exp(ta).
\end{equation}
There is a sign vector $s\in\{-1,1\}^N$ such that
\[
w=T(s\odot v),\qquad F(w)=v,\qquad
|w_j-1|<\varepsilon+a\quad\hbox{for every }j,
\qquad \sum_jw_j^2=N.
\]
If $\varepsilon+a<1$, then $w$ is strictly positive.
\end{lemma}

\begin{proof}
Put $c=1-\varepsilon$. Since $v_i\ge c>0$, assign the two signs the weights
\begin{equation}\label{eq:weights}
\Pr(s_i=1)=\frac{v_i+c}{2v_i},\qquad
\Pr(s_i=-1)=\frac{v_i-c}{2v_i}.
\end{equation}
These weights are nonnegative and sum to one. Use their finite product
distribution. For each coordinate,
\[
\E(s_iv_i)=c,\qquad
\E(s_iv_i-c)^2=v_i^2-c^2.
\]
All expectations here can be read simply as explicit finite weighted sums;
no measure-theoretic existence result is required.

For a fixed row $j$, set
\[
X_i=T_{ji}(s_iv_i-c),\qquad Z_j=\sum_iX_i=w_j-c.
\]
The row-sum identity in \eqref{eq:flat} gives the last equality. The variables
are centered, and flatness gives
\begin{equation}\label{eq:variance}
\sum_i\E X_i^2
=\frac1{d^2}\sum_i(v_i^2-c^2)=1-c^2\le2\varepsilon.
\end{equation}
Moreover,
\[
|X_i|\le\frac{v_i+c}{d}\le\frac2d.
\]
Thus $|tX_i|\le1$ by \eqref{eq:gates}. The elementary inequality
$e^z\le1+z+z^2$ for $|z|\le1$ implies
\[
\E e^{\pm tX_i}\le1+t^2\E X_i^2\le e^{t^2\E X_i^2}.
\]
Factoring the finite product distribution and using \eqref{eq:variance},
\begin{equation}\label{eq:mgf}
\E e^{tZ_j}\le e^{2\varepsilon t^2},\qquad
\E e^{-tZ_j}\le e^{2\varepsilon t^2}.
\end{equation}

Consider the nonnegative potential
\[
\Phi(s)=\sum_{j=1}^N\big(e^{tZ_j}+e^{-tZ_j}\big).
\]
Its weighted average is at most $2Ne^{2\varepsilon t^2}<e^{ta}$.
Hence at least one sign vector has $\Phi(s)<e^{ta}$. Every summand is
nonnegative, so $e^{tZ_j}<e^{ta}$ and $e^{-tZ_j}<e^{ta}$ for all $j$.
Since $t>0$, this gives $|Z_j|<a$, and consequently
$|w_j-1|<a+\varepsilon$.

Finally, $Tw=s\odot v$ by involutivity. Since $v>0$,
$F(w)=|s\odot v|=v$. Orthogonality and $s_i^2=1$ preserve the squared
Euclidean norm exactly.
\end{proof}

The small variance in \eqref{eq:variance} is essential to this argument.
A bound that uses only the ranges of the signs would discard the gain
obtained from $v$ being close to $\one$. The Lean proof establishes the
two-point moments, product factorization, exponential estimates, and
existence of a successful sign vector directly. It does not invoke an
unproved concentration inequality or a numerical random search.

\section{Explicit parameters and arbitrarily long positive orbits}
Fix $q\ge1$ and put
\begin{equation}\label{eq:params}
p=4^q,\qquad h=16p=2^{2q+4},\qquad d=h^p=2^{(2q+4)4^q},\qquad N=d^2.
\end{equation}
For $k\ge0$, define
\begin{equation}\label{eq:errors}
\varepsilon_k=\frac2{h^{4^k}},\qquad
a_k=\frac1{h^{4^k}},\qquad
t_k=\frac{(h^{4^k})^3}{8}.
\end{equation}
The subscript $k$ records the number of remaining error contractions, so
inverse steps use $k=q-1,q-2,\ldots,0$.

\begin{lemma}[All concentration gates hold]\label{lem:parameters}
For $0\le k<q$, Lemma~\ref{lem:inverse} applies with
$\varepsilon=\varepsilon_{k+1}$, $a=a_k$, and $t=t_k$.
Moreover,
\[
\varepsilon_{k+1}+a_k\le\varepsilon_k\le\frac18.
\]
\end{lemma}

\begin{proof}
Write $b=h^{4^k}$. Then $\varepsilon_{k+1}=2/b^4$,
$a_k=1/b$, and $t_k=b^3/8$. Since $b\ge h\ge16$,
$2/b^4+1/b\le2/b\le1/8$. Also $3\cdot4^k\le4^q$, so
$2t_k/d=b^3/(4d)\le1$.

The exponent identities are
\[
2\varepsilon_{k+1}t_k^2=\frac{b^2}{16},\qquad
t_ka_k=\frac{b^2}{8}.
\]
It remains to prove $2d^2<e^{b^2/16}$. Using $q\le4^q=p$,
$p\ge1$, and $\log2<1$,
\begin{align*}
\log(2d^2)
&=(1+4pq+8p)\log2\\
&<1+4pq+8p\le13p^2<16p^2
=\frac{h^2}{16}\le\frac{b^2}{16}.
\end{align*}
Multiplying the resulting strict inequality by $e^{b^2/16}$ proves
$2d^2e^{2\varepsilon_{k+1}t_k^2}<e^{t_ka_k}$.
\end{proof}

Choose any Bell label $j$ and define
\begin{equation}\label{eq:initial}
v_0=\one,\qquad v_1=T(\one-2e_j)=\one-2Te_j.
\end{equation}
Then
\[
F(v_1)=v_0,\quad Tv_1=\one-2e_j,\quad
\|v_1\|_2^2=N,\quad |v_1(i)-1|=\frac2d=\varepsilon_q.
\]
In particular $v_1>0$ and its partial-transpose coefficient at $j$ is $-1$.

Starting from $v_1$, apply the inverse lemma $q$ times, with the gates
provided by Lemma~\ref{lem:parameters}. This produces
\begin{equation}\label{eq:orbit}
v_0,v_1,\ldots,v_{q+1},\qquad
F(v_{i+1})=v_i,
\end{equation}
such that every coordinate of every $v_i$ is at least $7/8$, and every
$\|v_i\|_2^2=N$. The formal orbit interface retains the weaker bound $1/2$,
which is enough for the subsequent quantitative argument.

Every nonterminal reverse stage is also non-PPT. Indeed, if $Tv_i\ge0$
and $F(v_i)=v_{i-1}$, then $Tv_i=v_{i-1}$ and $Tv_{i-1}=v_i\ge0$.
Induction would contradict the known negative coordinate of $Tv_1$.
Only the defect at $v_1$ is needed by the central Lean theorem.

To obtain a definite finite specification, one may select at each step the
first sign vector, in a fixed order, satisfying the rational coordinate
bound for the next error level. Existence is proved by the preceding
argument. This selection is finite but astronomically expensive at the
displayed dimensions. The proof artifacts check the general existence
argument without expanding these state tables.

\section{Why full support forces an upper-level gap}
Reverse the constructed orbit:
\[
A_i=v_{q+1-i}\quad(0\le i\le q+1).
\]
Then
\begin{equation}\label{eq:forward}
A_{i+1}=F(A_i),\quad A_{q+1}=\one,\quad
A_i\ge\tfrac12\one,\quad \|A_0\|_2^2=N,\quad (TA_q)_j=-1.
\end{equation}
The greedy chain $A_0,A_1,\ldots,A_{q+1}$ is feasible at upper level $q+1$
and has terminal trace $N$.

We prove a strict lower bound on \emph{every} level-$q$ chain. This avoids
any need to infer a strict infimum gap merely from nonattainment of equality.

\begin{lemma}[Quantitative obstruction]\label{lem:gap}
If $s_0=A_0$, $s_{i+1}\ge|Ts_i|$ for $0\le i<q$, and $Ts_q\ge0$, then
\begin{equation}\label{eq:coarsegap}
\sum_js_q(j)\ge N+\delta_q,\qquad
\delta_q=\frac1{2(d+1)^{q+1}}>0.
\end{equation}
\end{lemma}

\begin{proof}
For each stage define $D_i=\diag(\operatorname{sign}(TA_i))$, choosing
sign $+1$ at zero. Thus
\[
A_{i+1}=D_iTA_i,\qquad TD_iA_{i+1}=A_i.
\]
Extend the candidate chain by $s_{q+1}=Ts_q$. Since $Ts_q\ge0$, this
extension satisfies the final sandwich as well and preserves total trace.
The active slack
\[
r_i=s_{i+1}-D_iTs_i
\]
is coordinatewise nonnegative. Orthogonality gives the exact identity
\begin{equation}\label{eq:slack}
A_{i+1}\cdot s_{i+1}-A_i\cdot s_i=A_{i+1}\cdot r_i\ge0.
\end{equation}
This underlying telescoping identity is the mechanism already used in
Brinster et al., Appendix C, Lemma 7~\cite{Brinster}. The present argument
uses an orbit with a uniform positive coordinate margin.

Put $\Delta=\sum_js_q(j)-N$. Telescoping \eqref{eq:slack} from $0$ to $q$
shows $\Delta\ge0$. Each individual increment is at most $\Delta$.
Since $A_{i+1}\ge\tfrac12\one$, for every coordinate and $i<q$,
\[
0\le r_i(j)\le2\Delta.
\]
Let $e_i=s_i-A_i$. Then $e_0=0$ and
\begin{equation}\label{eq:errorrec}
e_{i+1}=D_iTe_i+r_i.
\end{equation}
Flatness gives $\|Tz\|_\infty\le d\|z\|_\infty$. A direct induction in
\eqref{eq:errorrec} yields
\[
\|e_i\|_\infty\le2\Delta(d+1)^i\qquad(0\le i\le q).
\]
At the selected terminal coordinate, $Ts_q\ge0$ and $(TA_q)_j=-1$, hence
\[
1\le|(Te_q)_j|\le d\|e_q\|_\infty
\le2\Delta d(d+1)^q\le2\Delta(d+1)^{q+1}.
\]
This proves \eqref{eq:coarsegap}.
\end{proof}

The crucial distinction is between an arbitrarily deep absolute-transpose
orbit and an optimized SDP gap. An orbit with zero coordinates does not
control all active slacks. The positive margin in \eqref{eq:forward}
provides exactly that control. The proof does not assume a greedy optimizer.

A Euclidean-norm refinement gives a stronger conventional estimate, but it
is unnecessary for the central result. The theorem stated here and the verification
discussion use the conservative bound \eqref{eq:coarsegap}, which is
the one implemented in Lean.

\section{Normalization and proof of the main theorem}
Let $x=A_0=v_{q+1}$ and $\tau=\sum_jx_j$. Since $x_j\ge1/2$, we have
$\tau>0$. Define
\begin{equation}\label{eq:state}
\rho_q=\syn(x/\tau).
\end{equation}
All Bell coefficients are strictly positive and sum to one, so $\rho_q$
is a full-rank density matrix on the promised bipartite space.

The greedy chain synthesizes a feasible upper-level $(q+1)$ witness with
trace $N/\tau$. Thus
\[
\kappa_{q+1}(\rho_q)\le N/\tau.
\]
For any complex-matrix level-$q$ feasible chain, extract its Bell
expectations using \eqref{eq:expect} and multiply the coefficients by $\tau$.
This is a feasible raw chain with input $x$. Lemma~\ref{lem:gap} gives
\[
\operatorname{Re}\Tr S_{q-1}\ge\frac{N+\delta_q}{\tau}.
\]
Taking the infimum gives
\begin{equation}\label{eq:normalizedgap}
\kappa_q(\rho_q)\ge\frac{N+\delta_q}{\tau}
>\frac N\tau\ge\kappa_{q+1}(\rho_q),
\end{equation}
proving Theorem~\ref{thm:main}.

The formal proof establishes that the sets of feasible objective values are
nonempty and bounded below before applying conditional-infimum lemmas.
Nonemptiness follows, for example, by taking the first auxiliary coefficient
vector to be a sufficiently large constant and keeping that constant for
all subsequent stages. Its transform is the same nonnegative constant.
The terminal PPT condition implies a nonnegative trace, giving a lower bound.

The formalization therefore does not hide an empty optimization problem, an
assumed optimizer, an unproved duality theorem, or an artificial restriction
on the auxiliary matrices. The sandwich constraints also prove that the
real-trace formal objective agrees with the usual Hermitian SDP objective.

\section{Novelty assessment and significance}
\subsection{What was already known}
Lami, Mele, and Regula introduced converging upper and lower SDP hierarchies
and proved their relation to exact PPT cost~\cite{LMR}. Their Supplemental
Material, Conjecture S33, proposed finite collapse at a specific lower level,
with a corresponding upper-level formula. The same subsection asks the
analogous finite-collapse question for the upper hierarchy.

Brinster et al. construct arbitrary-depth separations for the lower $\chi$
hierarchy~\cite{Brinster}. Section IV, immediately after Eqs.\ (47)--(48),
explicitly observes that their states do not show upper $\kappa$ noncollapse
and leaves a finite upper formula open. Their Appendix C supplies the Bell
transform formalism, multinegative examples, and the telescoping dual argument.
These components receive prior credit; they are not claimed as discoveries
of the present work.

Earlier unpublished computations in the same research process had verified
only $\kappa_2>\kappa_3$ and $\kappa_3>\kappa_4$ on explicit states, with related
tensor nonmultiplicativity, and found one further finite example without a
complete Lean proof. None of those finite examples establishes the quantifier
over every $q$ proved here, and the present proof does not use them.

\subsection{The contribution assessed here}
The new mechanism is a positive inverse-orbit construction using biased
signs and a small-variance finite exponential potential. It supplies exactly
the full-support property that makes the upper optimization quantitatively
rigid. Together, these steps yield a strict gap at every adjacent finite
upper level for actual quantum states.

This resolves the specific remaining universal upper-hierarchy question in
the located primary literature. It is a structural result rather than a
small change in a constant or one additional numerical level. It explains
why optimized upper bounds can require arbitrarily large depth and removes
every fixed upper level as a universal exact-cost candidate.

The significance is bounded by the theorem's scope. The proof gives a very
large dimension bound and a small guaranteed normalized gap. It does not
improve the published approximation algorithm, optimize witness size, show
unbounded depth at a fixed dimension, or establish a computational-hardness
result. It leaves open other possible exact formulas. A specialist paper
on this result would need normal external review; no venue acceptance or
independent expert endorsement is asserted.

\subsection{Search coverage and qualified priority conclusion}
The audit refreshed the primary arXiv records and full texts, searched the
upper/lower terminology and normalized-body formulations, and checked
finite collapse, multinegativity, full-rank and dense examples, inverse
absolute partial transposition, positive signed-Hadamard orbits, and related
entanglement-cost additivity formulations. The normalized upper bodies
$\widetilde{\mathrm{PPT}}_q=\{X\ge0:\kappa_q(X)\le1\}$ were distinguished
from the lower bodies used in earlier contraction results.

Repository checks included the index and priority disclosures of this
collection,
\href{https://github.com/OpenAI/math}{OpenAI/math}, and the public
\href{https://github.com/the-omega-institute/trureturing}{trureturing}
problem/proof registry. The relevant quantum subtrees in the registry were
compared by tree identity with the preceding detailed audit; issue and PR
queries were refreshed. The located first-level $2\times3$ problem
\href{https://github.com/the-omega-institute/trureturing/issues/17758}{\#17758}
concerns a distinct $\kappa_1$ versus logarithmic-negativity question and
does not supply the present arbitrary-depth upper theorem.

A related September 2026 doctoral-thesis abstract was located in the Scuola
Normale Superiore repository; its full PDF returned an access error and was
not audited. Targeted MathOverflow and mathematical-database searches did
not locate a matching upper theorem. These access and indexing limitations
are retained in the novelty record.

The fresh Brinster arXiv record exposed only version 1, submitted
17 September 2026, when checked. No inspected source contained the
positive inverse-orbit result or all-depth upper separation proved here.
The search records of that audit are not part of this release.

\textbf{Priority assessment:} the theorem has a defensible, specifically
supported originality claim relative to the located literature. Historical
priority is still qualified. Search cannot exclude unindexed theses,
unpublished manuscripts, private ongoing work, or a proof filed under
unexpected terminology. Absence of a search hit is not treated as proof of
novelty. Lean verifies the mathematical implication, not historical priority.

\section{Lean correspondence and reproducibility}\label{sec:lean}
\subsection{What is formalized}
The central statement ranges over every natural $q\ge1$ and produces a
normalized complex bipartite density matrix. The proof includes the finite
sign existence lemma, all explicit parameter inequalities, the full finite
induction, the concrete Bell partial-transpose instance, the quantitative
coefficient obstruction, and its transfer to arbitrary complex-matrix SDP
chains. It is not a proof merely of a finite table, a numerical residual,
or an abstract involution without a physical instance.

The proof is organized as follows.
\begin{center}\small
\begin{tabular}{@{}p{.31\textwidth}p{.62\textwidth}@{}}
\toprule
Module & Mathematical role\\\midrule
\texttt{FoldSystem} & Flat symmetric involution; norm, trace, and dual-pairing identities.\\
\texttt{SignConcentration} & Two-point weights and moments; finite product exponential bound.\\
\texttt{InverseSign} & Simultaneous coordinate control and an exact positive inverse step.\\
\texttt{ParameterBounds} & Explicit dimension, error, tilt, and strict exponential gates.\\
\texttt{ReverseOrbit} & Initial defect and the complete finite reverse-orbit induction.\\
\texttt{BellFoldSystem} & Actual product-Bell partial transposition satisfies every abstract assumption.\\
\texttt{ChainGap} & Positive lower bound for every shorter feasible coefficient chain.\\
\texttt{ArbitraryHierarchy} & Original arbitrary-depth complex SDP, nonempty infima, normalization, and matrix gap.\\
\texttt{BellPositivity} & Strictly positive Bell weights imply positive definiteness.\\
\texttt{AllDepth} & Final all-depth physical theorem and no universal finite collapse.\\
\bottomrule
\end{tabular}
\end{center}

Earlier small structural modules provide Bell projectors, synthesis,
expectations, and PSD facts. The earlier research's large finite data tables
and optimization scripts are not dependencies of the new theorem.

The use of \texttt{noncomputable} marks ordinary real-valued definitions and
existential choices. It does not add axioms, bypass the kernel, or replace
proofs with a solver result. Every sign-existence claim is proved. The
project uses neither \texttt{sorry}/\texttt{admit}, custom axioms, nor a native
decision oracle. Only the usual Lean foundations \texttt{propext},
\texttt{Classical.choice}, and \texttt{Quot.sound} are allowed by the audit.

\subsection{Pinned clean replay}
The release pins Lean 4.19.0, commit \texttt{6caaee842e94}, and mathlib at
\begin{center}\small\texttt{c44e0c8ee63ca166450922a373c7409c5d26b00b}.\end{center}
Every transitive dependency is pinned by full git commit. The verifier
rejects dirty dependency checkouts, changed frozen sources, unlisted project
imports, prohibited proof shortcuts, unexpected axioms, and incomplete
audit output.

The replay freshly compiles 17 mathematical modules and
the generated axiom audit, checking 256 public source declarations.
Each compiler invocation checks the kernel at trust zero and disables
automatic implicit parameters. The exact flags are
\begin{center}\small\verb|-j1 --trust=0 -DwarningAsError=true -DautoImplicit=false|\end{center}
The replay recorded with this release used a Mathlib checkout at the
pinned revision whose modules had been built from source on the same machine;
it took about seven minutes with two parallel jobs.

\begin{minipage}{\textwidth}
From the \texttt{lean/} directory of this package, after preparing the exact
Lean release and dependencies, run:
\begin{verbatim}
lake exe cache get
python3 verify_release.py \
  --lean-home /absolute/path/to/lean-4.19.0-linux \
  --mathlib .lake/packages/mathlib \
  --packages-dir .lake/packages \
  --out /absolute/path/to/new-empty-verification-directory \
  --jobs 2
\end{verbatim}
\end{minipage}

Do not run \texttt{lake update}, which changes the reviewed dependency
manifest. A pristine standalone mathlib checkout at the same pin may also
be supplied; \texttt{REPRODUCE.txt} gives that command.

Success requires a receipt with status \texttt{PASS}, zero exit status for
every module, and an axiom inventory matching the complete reviewed list.
The receipt records exact commands, source/object/log hashes, dependency
revisions, and timing. The receipt and logs of the replay for this release are
in the package's \texttt{verification/} directory, and the continuous
integration workflow repeats the replay on every change.

\section{Independent scrutiny and remaining limitations}
In the research process, AI reviewers challenged the proof on four points: whether the inverse
sign distribution is genuinely defined and normalized; whether its variance
and simultaneous bounds have the correct dimension factors; whether full
support controls an optimized upper chain; and whether the formal statement
is the actual complex-matrix SDP. Each issue is addressed by explicit
lemmas, and the final theorem discharges the conditional orbit hypotheses.

The development explored finite reverse-mask searches, dual exposed faces,
and fixed-dimension nonpolyhedral possibilities before the general weighted
sign argument. Those explorations did not establish the all-depth theorem
and are not used as proof evidence. The old flag-padding, face-extraction,
and greedy-depth shortcuts were not revived as assumptions. The successful
proof proceeds through the complete inverse lemma and quantitative slack
estimate given above.

The final artifacts establish mathematical correctness subject to the
standard trust in Lean's kernel and the pinned library foundations.
They do not replace specialist peer review of the theorem's presentation,
priority, or practical interpretation. The main unresolved quantitative
questions are much smaller witness dimensions, useful normalized gaps,
and simple successful sign patterns. These are future questions rather
than qualifications to the proved all-depth existence theorem.

\section*{Research process and AI disclosure}
This is fully autonomous agentic research. Dev D. Goyal is responsible for publication. No expert validation or external peer review is claimed, and the novelty assessment is provisional.

The theorem, its proof, the Lean formalization, the literature searches and an earlier version of this note were produced by autonomous AI research agents in an earlier session. The records of that session do not name the AI system used.

Opus 5.5 by Anthropic used with Claude Code assessed the package for this release. It checked the definition of $\kappa_q$ in the Lean sources against the hierarchy stated in Section~1, including the partial transpose on the second system, the terminal PPT condition and the unrestricted complex auxiliaries; rederived by hand the inverse-sign lemma, the parameter inequalities, the initial defect and the quantitative chain gap; ran the fresh Lean replay described above; and revised this manuscript. The primary arXiv records could not be opened from its environment, so the comparison with the cited literature relies on the earlier session's audit. AI checks are not human peer review. Lean verification establishes only the formal statements described in Section~\ref{sec:lean}.

\begingroup
\interlinepenalty=10000
\small
\begin{thebibliography}{9}
\setlength{\itemsep}{3pt}\setlength{\parsep}{0pt}
\bibitem{LMR}
L. Lami, F. A. Mele, and B. Regula.
\emph{Computable entanglement cost under positive partial transpose operations}.
Physical Review Letters \textbf{134}, 090202 (2025).
\href{https://arxiv.org/abs/2405.09613v2}{arXiv:2405.09613v2}, 7 March 2025.

\bibitem{Brinster}
R. Brinster, T. V. Kondra, H. Kampermann, D. Bru\ss, and N. Wyderka.
\emph{Multinegativity and single-letter formulas for asymptotic entanglement}.
\href{https://arxiv.org/abs/2609.20698v1}{arXiv:2609.20698v1}, 17 September 2026.

\bibitem{WJZ}
X. Wang, M. Jing, and C. Zhu.
\emph{Computable and Faithful Lower Bound on Entanglement Cost}.
\href{https://arxiv.org/abs/2311.10649v4}{arXiv:2311.10649v4}, 7 May 2025.

\bibitem{Lean}
The Lean and mathlib communities.
\emph{Lean 4.19.0 and mathlib}, pinned source and dependency manifest supplied
with this release. Project proof sources, the replay receipt and the full
axiom inventory are included with this manuscript's package.
\end{thebibliography}
\endgroup
\end{document}
